\section{The Shipped Detector: Threshold, Gate, Debounce}
\label{sec:detector}

Everything in the preceding chapters converges on a single, deliberately small
decision rule. Out of all the spectral machinery, Cadence ships exactly one
on-board verdict per control window, and that verdict drives the OLED, the
NeoPixel colour, and the web dashboard. This is the \emph{glass box}: a rule a
clinician can read in one line and trust. The whole detector is

\begin{equation}
\textsf{FREEZE} \iff
\underbrace{(\text{movement energy} > \text{still\_floor})}_{\text{the leg is actually moving}}
\ \wedge\
\underbrace{(\FI > \tau)}_{\text{moving like a freeze}},
\quad\text{then \emph{debounced}.}
\label{eq:rule}
\end{equation}

\noindent Three ingredients, in this order: a \textbf{gate} (is the leg moving at
all?), a \textbf{threshold} on the Freeze Index $\FI$ (is the movement the
3--8\,Hz tremble of a freeze rather than the swing of a walk?), and a
\textbf{debounce} (has the new verdict held for long enough to be believed?). The
deployed trip line is $\tau = 1.815$. The rest of this section explains why each
piece exists, why $\tau$ is set below the textbook optimum, and how the three
combine into the STILL\,/\,WALKING\,/\,FREEZE state machine.

\begin{keyidea}
The detector is not a black box that emits a probability. It is a short Boolean
sentence over two physically meaningful quantities --- how \emph{hard} the leg is
moving (movement energy) and how \emph{freeze-like} that motion is ($\FI$). Both
can be read off the live dashboard, so every verdict is auditable after the fact.
\end{keyidea}

\subsection{Why a gate comes first: the still-leg problem}

Recall that the Freeze Index is a \emph{ratio} of band powers,
$\FI = P_{3\text{--}8} / P_{0.5\text{--}3}$. Ratios are scale-invariant, which is
exactly what we want when comparing a freeze to a walk --- but it is also their
weakness when there is almost no signal at all. Stand perfectly still and both
band powers collapse towards the sensor's noise floor. Dividing one small noisy
number by another small noisy number can produce \emph{any} value, including a
large one. A motionless leg can therefore manufacture a spurious high $\FI$ out of
pure electronic hash, and a $\FI$-only detector would cry ``freeze'' at someone
sitting in a chair.

The fix is a second quantity that is deliberately \emph{not} scale-invariant:

\begin{equation}
\text{movement energy} \;=\; P_{0.5\text{--}3}\;+\;P_{3\text{--}8},
\label{eq:energy}
\end{equation}

\noindent the total band power across the locomotor and freeze bands of the
mean-removed magnitude. Unlike the ratio, this number reports \emph{how much the
leg is moving} in absolute terms. A still leg has near-zero movement energy; a
walking or freezing leg has plenty. The gate is a simple floor on this quantity:
if movement energy is below \code{still\_floor}, the device declares STILL and
never even consults the $\FI$.

\begin{mathsbox}
$\FI$ and movement energy are built from the \emph{same two band powers}, combined
two different ways:
\[
\FI=\frac{P_{3\text{--}8}}{P_{0.5\text{--}3}}\ \ (\text{a \emph{ratio} --- shape}),
\qquad
\text{energy}=P_{0.5\text{--}3}+P_{3\text{--}8}\ \ (\text{a \emph{sum} --- size}).
\]
The ratio answers \emph{``what kind of motion?''} and the sum answers \emph{``is
there any motion?''} You need both: shape without size is the still-leg trap;
size without shape cannot tell a walk from a freeze.
\end{mathsbox}

\code{still\_floor} is not a magic constant. It is \textbf{auto-calibrated at
boot}: when the firmware starts, it samples the resting movement energy of the
ankle for a moment and sets the floor just above that quiet baseline. This adapts
the gate to the individual sleeve fit, the mounting tightness, and the particular
board's sensor offsets, so ``still'' means \emph{this} wearer's still, not a
factory guess.

\begin{intuition}
Think of the gate as a bouncer at the door. Before anyone asks ``is this person
dancing or freezing?'' the bouncer first checks ``is this person moving at all?''
If the leg is parked, the question of \emph{how} it is moving is meaningless ---
and the noisy $\FI$ that the maths would hand back is simply ignored. On the
dashboard this is the \textbf{Movement Energy} panel's MOVING\,/\,STILL badge, and
when it reads STILL the Freeze Margin meter shows ``\,---\,/\,at rest.''
\end{intuition}

\subsection{The threshold: 1.815 versus the textbook optimum 2.10}

Once the gate confirms the leg is moving, the detector asks whether the motion is
freeze-shaped by comparing $\FI$ against the trip line $\tau$. Choosing $\tau$ is a
classic sensitivity/specificity trade-off, and Cadence makes a deliberate,
clinically-motivated choice that differs from the value a naive optimiser would
pick.

\begin{figure}[H]
\centering
\includegraphics[width=0.86\linewidth]{operating_point}
\caption{How an operating point is chosen --- a threshold sweep. As the trip line
moves, \textcolor{freeze}{sensitivity} (freezes caught) falls while
\textcolor{navy}{specificity} (calm gait correctly left alone) rises; the dashed
curve is Youden's $J=\text{sensitivity}+\text{specificity}-1$, whose peak (star)
is the ``clean'' balanced optimum. Cadence faces the same trade-off and chooses to
sit \emph{left} of the balanced point, buying recall of freezes at the cost of a
few false alarms. (Curve shown for the interpretable model of
Sec.~\ref{sec:detector}'s companion study; the principle is identical for the
$\FI$ trip line.)}
\label{fig:opsweep}
\end{figure}

To pick $\tau$ for the shipped rule we ran the real worksheet capture,
\code{capture1.csv} (64\,Hz, 41.4\,s), through the exact windowing the firmware
uses and counted outcomes window-by-window. The capture yields 19 control
windows: 6 are genuine freeze windows and 13 are non-freeze (still or walking).
Sweeping the trip line gives the two candidate operating points in
Table~\ref{tab:sweep}.

\begin{table}[H]
\centering
\caption{Real threshold sweep on \code{capture1.csv} (19 windows: 6 freeze,
13 non-freeze). TP/FP/TN/FN counted per window. The shipped board uses
$\tau=1.815$; $\tau=2.10$ is the balanced (Youden) optimum.}
\label{tab:sweep}
\small
\begin{tabular}{l cccc cc l}
\toprule
Trip line $\tau$ & TP & FP & TN & FN & Sensitivity & Specificity & Note \\
\midrule
$\FI>1.815$ & 5 & 1 & 12 & 1 & \textbf{0.83} & 0.92 & \emph{shipped} (sensitivity-biased) \\
$\FI>2.10$  & 5 & 0 & 13 & 1 & 0.83 & \textbf{1.00} & clean Youden optimum \\
\bottomrule
\end{tabular}
\end{table}

Read the table carefully, because the lesson is subtle. Both thresholds catch the
same 5 of 6 freezes (sensitivity $0.83$ either way) and both miss the same one
short freeze. The \emph{only} difference is a single window: at $\tau=1.815$ one
calm window is wrongly flagged (FP\,$=1$, specificity $0.92$), whereas at
$\tau=2.10$ that window is correctly left alone (FP\,$=0$, specificity $1.00$). So
the higher threshold is strictly cleaner on \emph{this} capture --- it is the
Youden-optimal point. Why does the board ship the lower, ``worse'' number?

\begin{takeaway}
A missed freeze can mean a fall; a false alarm just lights an LED that nobody acts
on. The costs are wildly asymmetric, so Cadence biases toward \textbf{sensitivity}
and ships $\tau=1.815$. On this capture both thresholds happen to tie on
sensitivity, but $1.815$ leaves head-room: it keeps catching freezes that sit just
\emph{below} $2.10$ on unseen wearers, where the clean optimum would start
dropping them.
\end{takeaway}

The point is that $2.10$ is optimal \emph{for capture1}, and over-fitting a single
short capture is exactly the trap to avoid. The trade-off the firmware actually
cares about is the one on the \emph{next} wearer, whose freezes we have not seen.
Pushing the trip line down to $1.815$ is cheap insurance: it costs one false alarm
on this recording (a harmless flicker on the display) and buys a margin of safety
against the freezes that a tighter line would let slip. The asymmetry of harm,
not the arithmetic of one CSV, sets the threshold.

\begin{keyidea}
There are \emph{two} different ``right'' answers. The \textbf{Youden} answer
($2.10$) maximises balanced accuracy on the data in front of you. The
\textbf{clinical} answer ($1.815$) minimises the expected \emph{harm}, given that
a false negative (missed freeze $\to$ possible fall) costs far more than a false
positive (a brief, ignorable LED). Cadence is a safety monitor, so it ships the
clinical answer.
\end{keyidea}

\noindent A note on provenance: a library default constant elsewhere in the code
base reads $2.0$, but the \emph{deployed} value baked into the monitor firmware is
$1.815$. The two should not be confused --- only $1.815$ is the shipped trip line.

\subsection{Seeing the rule on real data}

Figure~\ref{fig:cap1} makes the whole rule concrete on \code{capture1}. The top
panel is the raw accelerometer magnitude $\norm{a}$ in milli-g; the bottom panel is
the per-window Freeze Index. The capture is staged in four phases: still, walking,
freeze, then recovery.

\begin{figure}[H]
\centering
\includegraphics[width=0.92\linewidth]{capture1_trace}
\caption{The detector's two inputs on \code{capture1}. \emph{Top:} accelerometer
magnitude $\norm{a}$ (mg) --- note the violent, dense oscillation during the
freeze phase, with heel-strike-scale peaks that would clip a $\pm2$\,g range.
\emph{Bottom:} the per-window Freeze Index, which sits near $1$ through still and
walking and erupts to a peak of $21.35$ during the freeze. The shipped trip line
$\tau=1.815$ cuts cleanly between the two regimes.}
\label{fig:cap1}
\end{figure}

The numbers match the picture. Per-phase $\FI$ on this capture is around $1.1$ at
rest (max $1.71$) and around $1.2$ while walking (max $2.08$); during the freeze it
jumps to a mean of $6.89$, median $4.45$, and peak $21.35$. The trip line at
$1.815$ sits comfortably in the empty gap above the walking distribution and well
below the freeze surge --- which is precisely why a simple horizontal threshold
works at all. Note too that the still and walking maxima ($1.71$ and $2.08$)
straddle $1.815$: that lone window above the line during normal activity is the
single false positive in Table~\ref{tab:sweep}, and is exactly what the gate and
debounce exist to tame.

\subsection{Debounce: don't trust a one-window verdict}

A 256-sample window with a 128-sample hop produces one fresh verdict every two
seconds. Real motion is noisy, and a single unlucky window --- one heel-strike
landing awkwardly, one transient bump --- can momentarily push $\FI$ across the
line. If the device committed every window instantly, the dashboard would flicker
and a one-off jolt could read as a freeze.

Cadence therefore \textbf{debounces}: a \emph{new} state must persist for
\textbf{two consecutive} control windows before it is committed. This is
\emph{assert-after-2} for entering a state and \emph{release-after-2} for leaving
one. A single anomalous window can never flip the verdict; it takes a sustained,
two-window-deep change ($\geq$ about four seconds of agreeing evidence given the
overlap) to move the committed state. Because a genuine freeze lasts seconds, real
episodes clear this bar easily, while isolated spikes do not.

\begin{intuition}
Debounce is the same trick as a mechanical light switch that must be held, not
brushed. One stray window is a brush; the switch ignores it. Two windows in a row
saying the same thing is a deliberate press, and only then does the verdict
change. The cost is a short, bounded latency --- but a freeze that lasts seconds
has time to spare, and a fall-risk monitor would rather be right than instant.
\end{intuition}

\subsection{Putting it together: the state machine}

The gate, the threshold, and the debounce compose into a three-state machine that
emits STILL, WALKING, or FREEZE. Figure~\ref{fig:tree} draws the live decision
path for a single control window.

\begin{figure}[H]
\centering
\resizebox{0.95\textwidth}{!}{%
\begin{tikzpicture}[
  font=\sffamily,
  >=Stealth,
  node distance=14mm,
  deci/.style={diamond, aspect=1.7, draw=navy, fill=palenavy, line width=0.7pt,
               inner sep=1pt, align=center, minimum width=34mm},
  state/.style={rounded corners=3pt, draw, line width=0.8pt, align=center,
                minimum width=26mm, minimum height=11mm, font=\bfseries},
  still/.style={state, draw=mute,   fill=mist},
  walk/.style={state,  draw=grn,    fill=paleteal},
  frz/.style={state,   draw=freeze, fill=palefreeze},
  lbl/.style={font=\footnotesize\sffamily, fill=white, inner sep=1.5pt},
  flow/.style={->, line width=0.8pt, draw=ink}
]

% input
\node[rounded corners=3pt, draw=teal, fill=paleteal, line width=0.8pt,
      align=center, minimum height=11mm, minimum width=30mm] (in)
      {New 4\,s window\\[1pt]\footnotesize\mdseries energy, $\FI$};

% gate
\node[deci, right=20mm of in] (gate) {energy $>$\\ still\_floor\,?};

% FI test
\node[deci, right=24mm of gate] (fi) {$\FI > 1.815$\,?};

% debounce (shared) — placed to the right of FI
\node[deci, right=24mm of fi] (deb) {same state\\ 2 windows\\ in a row\,?};

% states column on the far right
\node[frz,   right=22mm of deb] (freeze) {FREEZE};
\node[walk, below=16mm of freeze] (walking) {WALKING};
\node[still, above=16mm of freeze] (stillS) {STILL};

% flows
\draw[flow] (in) -- (gate);
\draw[flow] (gate) -- node[lbl,above]{moving} (fi);
\draw[flow] (fi)   -- node[lbl,above]{yes} (deb);

% gate NO -> still candidate (down then to debounce as STILL)
\coordinate (gbot) at ($(gate.south)+(0,-20mm)$);
\draw[flow] (gate.south) -- node[lbl,right=1pt]{still} (gbot)
            -| (deb.south);
% FI NO -> walking candidate
\coordinate (fbot) at ($(fi.south)+(0,-11mm)$);
\draw[flow] (fi.south) -- node[lbl,right=1pt]{no}(fbot) -| ($(deb.south)+(-3mm,0)$);

% debounce YES -> commit (to whichever candidate state was carried)
\draw[flow] (deb.east) -- node[lbl,above]{commit}(freeze.west);
\draw[flow] (deb.north) -- node[lbl,left=1pt]{if STILL} (stillS.south);
\draw[flow] (deb.south) -- node[lbl,left=1pt]{if WALK} (walking.north);

% debounce NO -> hold (loop)
\draw[flow] (deb.north west) .. controls ($(deb.north west)+(-4mm,12mm)$) and
            ($(fi.north east)+(4mm,12mm)$) ..
            node[lbl,above]{no: hold last state} (fi.north east);

\end{tikzpicture}}
\caption{The shipped state machine for one control window. The \textbf{gate}
(energy vs.\ \code{still\_floor}) runs first; only if the leg is moving is the
\textbf{$\FI$ threshold} consulted. The resulting candidate state
(STILL\,/\,WALKING\,/\,FREEZE) is committed only after the \textbf{debounce}
confirms it has held for two consecutive windows; otherwise the previous committed
state is retained. Equivalently: \emph{STILL} if not moving; else \emph{FREEZE} if
$\FI>1.815$; else \emph{WALKING}.}
\label{fig:tree}
\end{figure}

In plain words, the state for each window is decided as:
\begin{enumerate}
\item \textbf{Gate.} If movement energy $<$ \code{still\_floor}, the candidate is
\textsf{STILL} and the $\FI$ is not even evaluated (it would be noise).
\item \textbf{Threshold.} Otherwise, if $\FI > 1.815$ the candidate is
\textsf{FREEZE}; if not, the leg is moving but not freeze-shaped, so the candidate
is \textsf{WALKING}.
\item \textbf{Debounce.} The candidate becomes the \emph{committed} state only once
it has held for two consecutive windows; until then, the last committed state
stands.
\end{enumerate}

This committed state is the single output of the device. It paints the NeoPixel
ring (state colours), prints to the OLED, and on the live dashboard drives the
headline \textbf{Gait Status} panel, while the \textbf{Freeze Margin} meter shows
the live $\FI$ as a percentage of the $1.815$ trip line so a user can watch a
freeze build before the line is crossed.

\begin{takeaway}
The whole shipped detector is three guards in series --- \textbf{gate} (is the leg
moving?), \textbf{threshold} (is it moving like a freeze?), \textbf{debounce} (has
it held?). Each one removes a specific failure mode: the gate kills the still-leg
$\FI$ blow-up, the sensitivity-biased threshold ($1.815$) refuses to miss a freeze,
and the two-window debounce refuses to be fooled by a single jolt. Simple,
inspectable, and tuned for the cost that actually matters: a missed freeze.
\end{takeaway}
